% Part: incompleteness
% Chapter: introduction
% Section: decidability

\documentclass[../../../include/open-logic-section]{subfiles}

\begin{document}

\olfileid{inc}{int}{dec}

\olsection{تصمیم‌ناپذیری و ناتمامیت}

اثبات گودل از قضایای ناتمامیت به حسابی‌سازی نحو نیاز دارد. اما حتی
بدون آن نیز می‌توانیم تنها با فرض اینکه نظریه‌ای همهٔ روابط را
!!{represents} و این روابط !!{decidable}ند، به نتایج خوبی برسیم. اثبات،
استدلالی قطری شبیه اثبات تصمیم‌ناپذیری مسئلهٔ توقف است.

\begin{thm}
اگر $\Gamma$ نظریه‌ای سازگار باشد که هر رابطه را !!{represents}
و آن رابطه !!{decidable} است، آنگاه $\Gamma$ !!{decidable} نیست.
\end{thm}

\begin{proof}
فرض کنید $\Gamma$ !!{decidable} باشد. نشان می‌دهیم اگر $\Gamma$ هر
رابطه را !!{represents} و آن رابطه !!{decidable} باشد، باید ناسازگار باشد.

ویژگی‌های !!^{decidable} (روابط یک‌موضعی) با !!{formula}sیی دارای یک
متغیر آزاد بازنمایی می‌شوند. بگذارید $!A_0(x)$، $!A_1(x)$،~\dots
برشمردنی محاسبه‌پذیر از همهٔ چنین !!{formula}sیی باشد. اکنون مجموعهٔ
زیر را در نظر بگیرید:
\[
D = \Setabs{n}{\Gamma \Proves \lnot !A_n(\num{n})}.
\]
مجموعهٔ~$D$ !!{decidable} است، زیرا برای آزمودن $n\in D$ نخست
$!A_n(x)$ و سپس $\lnot !A_n(\num{n})$ را محاسبه می‌کنیم. روشن است که
جانشانی ترم~$\num{n}$ به جای هر رخداد آزاد~$x$ در~$!A_n(x)$ و افزودن
$\lnot$ به آغاز~$!A_n(\num{n})$ کاری مکانیکی است. بنا بر فرض،
$\Gamma$ !!{decidable} است، پس می‌توانیم بیازماییم آیا
$\lnot !A_n(\num{n})\in\Gamma$. اگر هست، $n\in D$ و اگر نیست،
$n\notin D$. پس $D$ نیز !!{decidable} است.

چون $\Gamma$ همهٔ ویژگی‌ها را !!{represents} و این ویژگی‌ها
!!{decidable}ند،
$D$ را نیز !!{represents}. و !!{formula}sیی که $D$ را در
$\Gamma$ !!{represents}s همگی در میان
$!A_0(x)$، $!A_1(x)$،~\dots هستند. پس عددی چون~$d$ بگیرید که
$!A_d(x)$، $D$ را در~$\Gamma$ !!{represents}. اگر $d\notin D$، چون
$!A_d(x)$، $D$ را !!{represents}،
$\Gamma\Proves\lnot !A_d(\num{d})$. اما این یعنی $d$ شرط تعریف‌کنندهٔ
$D$ را برآورده می‌کند و بنابراین $d\in D$. این با $d\notin D$ تناقض
دارد. پس با برهان خلف، $d\in D$.

چون $d\in D$، بنا بر تعریف~$D$،
$\Gamma\Proves\lnot !A_d(\num{d})$. از سوی دیگر، چون $!A_d(x)$، $D$
را در~$\Gamma$ !!{represents}،
$\Gamma\Proves !A_d(\num{d})$. در نتیجه $\Gamma$ ناسازگار است.
\end{proof}

\begin{explain}
قضیهٔ پیشین نشان می‌دهد هیچ نظریهٔ سازگاری که همهٔ روابط را
!!{represents} و این روابط !!{decidable}ند، نمی‌تواند !!{decidable} باشد.
نشان خواهیم داد $\Th{Q}$ همهٔ روابط را !!{represents}s و این روابط !!{decidable}ند؛ یعنی همهٔ نظریه‌هایی که~$\Th{Q}$ را در بر دارند، مانند
$\Th{PA}$ و~$\Th{TA}$، نیز چنین‌اند و بنابراین !!{decidable}
نیستند. (چون همهٔ این نظریه‌ها در مدل استاندارد صادق‌اند، همگی
سازگارند.)

همچنین می‌توانیم با این نتیجه نسخه‌ای ضعیف از قضیهٔ نخست ناتمامیت به
دست آوریم. هر نظریه‌ای که !!{axiomatizable} و !!{complete} باشد
!!{decidable} است. پس نظریه‌های سازگار و !!{axiomatizable}ی که همهٔ
ویژگی‌ها را !!{represents}s و این ویژگی‌ها !!{decidable}ند، نمی‌توانند !!{complete}
باشند.
\end{explain}

\begin{thm}
اگر $\Gamma$ !!{axiomatizable} و !!{complete} باشد، !!{decidable} است.
\end{thm}

\begin{proof}
هر نظریهٔ ناسازگاری !!{decidable} است، زیرا نظریه‌های ناسازگار همهٔ
!!{sentence}s را در بر دارند. پس پاسخ پرسش «آیا $!A\in\Gamma$؟» همیشه
«بله» است و بنابراین تصمیم‌پذیر است.

اکنون فرض کنید $\Gamma$ سازگار و افزون بر آن !!{axiomatizable} و
!!{complete} باشد. چون $\Gamma$ !!{axiomatizable} است،
!!{computably enumerable} است. زیرا می‌توانیم همهٔ !!{derivation}sی درست
از اصول موضوعهٔ~$\Gamma$ را با تابعی محاسبه‌پذیر برشماریم. از هر
!!{derivation} درست می‌توانیم !!{sentence}ای را که !!{derive}s محاسبه کنیم؛
پس در مجموع تابعی محاسبه‌پذیر وجود دارد که همهٔ قضایای~$\Gamma$ را
برمی‌شمارد. !!{sentence}~$!A$ قضیهٔ~$\Gamma$ است اگر و تنها اگر
$\lnot !A$ قضیه نباشد، زیرا $\Gamma$ سازگار و !!{complete} است. بنابراین می‌توانیم عضویت
$!A$ در~$\Gamma$ را چنین تصمیم کنیم. همهٔ قضایای~$\Gamma$ را
برشمارید. هنگامی که $!A$ در فهرست ظاهر شود، می‌دانیم
$\Gamma\Proves !A$. هنگامی که $\lnot !A$ ظاهر شود، می‌دانیم
$\Gamma\Proves/ !A$. چون $\Gamma$ !!{complete} است، سرانجام یکی از این
دو حالت رخ می‌دهد و روش پاسخی تولید می‌کند.
\end{proof}

\begin{cor}
\ollabel{cor:incompleteness}
اگر $\Gamma$ سازگار و !!{axiomatizable} باشد و هر ویژگی را
!!{represents} و آن ویژگی !!{decidable} باشد، !!{complete} نیست.
\end{cor}

\begin{proof}
اگر $\Gamma$ !!{complete} بود، بنا بر قضیهٔ پیشین !!{decidable} می‌بود
(زیرا !!{axiomatizable} و سازگار است). اما چون هر ویژگی را
!!{represents} و آن ویژگی !!{decidable} است، بنا بر قضیهٔ نخست
!!{decidable} نیست.
\end{proof}

\begin{prob}
نشان دهید $\Th{TA}=\Setabs{!A}{\Sat{N}{!A}}$
!!{axiomatizable} نیست. می‌توانید فرض کنید $\Th{TA}$ همهٔ
ویژگی‌های تصمیم‌پذیر را بازنمایی می‌کند.
\end{prob}

وقتی ثابت کنیم مثلاً $\Th{Q}$ همهٔ ویژگی‌ها را !!{represents}
و این ویژگی‌ها !!{decidable}ند، نتیجهٔ فرعی به ما می‌گوید $\Th{Q}$ باید ناتمام
باشد. اما اثبات آن نمونه‌ای از !!{sentence}ای مستقل به دست نمی‌دهد؛ فقط
نشان می‌دهد چنین !!a{sentence}ای باید وجود داشته باشد. برای یافتن نمونه،
باید نحو را حسابی‌سازی و ایدهٔ اصلی اثبات گودل را دنبال کنیم. و البته
هنوز باید ادعای نخست را نشان دهیم: اینکه $\Th{Q}$ واقعاً همهٔ
ویژگی‌ها را !!{represents}s و این ویژگی‌ها !!{decidable}ند.

باید توجه کرد که هر نظریهٔ \emph{جالبی} ناتمام یا تصمیم‌ناپذیر نیست.
نظریه‌های بسیاری چنان قوی هستند که واقعیت‌های ریاضی جالبی را توصیف
کنند، اما شرایط نتیجهٔ گودل را برآورده نمی‌کنند. برای نمونه،
$\Th{Pres}=\Setabs{!A\in\Lang{L_{A^+}}}{\Sat{N}{!A}}$، یعنی مجموعهٔ
!!{sentence}sی صادق در مدل استاندارد از زبان حسابِ بدون~$\times$، هم
کامل و هم تصمیم‌پذیر است. این نظریه حساب پرسبرگر نام دارد و همهٔ
حقایق دربارهٔ اعداد طبیعی را که تنها با~$\Obj{0}$، $\prime$ و~$+$
صورت‌بندی‌پذیرند ثابت می‌کند.

\end{document}
